\newpage
\phantomsection
\label{prf:euclidean-distance-rn-metric}
\LRAProofFor{prop:euclidean-distance-rn-metric}

\begin{remark*}[Return]
\hyperref[prop:euclidean-distance-rn-metric]{Return to Proposition}
\end{remark*}

\begin{proposition*}[The Euclidean Distance on \(\mathbb{R}^n\) is a Metric]
The function \(d_2\) is a metric on \(\mathbb{R}^n\).
\end{proposition*}

\begin{proof}[Professional Standard Proof]
\LRAProofBodyStart
For \(x,y\in\mathbb{R}^n\),
\[
  d_2(x,y)=\sqrt{\sum_{i=1}^n (x_i-y_i)^2}\geq 0.
\]
This distance is \(0\) if and only if every square \((x_i-y_i)^2\) is \(0\),
which is equivalent to \(x_i=y_i\) for every \(i\), hence to \(x=y\).
Symmetry follows from \((x_i-y_i)^2=(y_i-x_i)^2\).

For the triangle inequality, set \(u=x-z\) and \(v=z-y\).  Then
\[
  d_2(x,y)=|u+v|,\qquad d_2(x,z)=|u|,\qquad d_2(z,y)=|v|.
\]
Using \hyperref[thm:cauchy-schwarz-inequality]{Cauchy--Schwarz},
\[
  |u+v|^2
  = |u|^2+2u\cdot v+|v|^2
  \leq |u|^2+2|u|\,|v|+|v|^2
  = (|u|+|v|)^2.
\]
Both sides are nonnegative, so \( |u+v|\leq |u|+|v| \).  Therefore
\[
  d_2(x,y)\leq d_2(x,z)+d_2(z,y).
\]
Thus \(d_2\) is a metric.
\end{proof}

\begin{proof}[Detailed Learning Proof]
\LRAProofBodyStart
The first three metric requirements are coordinate checks.  A sum of squares
is nonnegative, and it vanishes exactly when every coordinate difference
vanishes.  Symmetry holds because changing \(x_i-y_i\) to \(y_i-x_i\) only
changes the sign before squaring.

The triangle inequality is the vector-length statement
\[
  |u+v|\leq |u|+|v|.
\]
After squaring, it becomes
\[
  |u+v|^2\leq (|u|+|v|)^2.
\]
Expanding the left side gives \(|u|^2+2u\cdot v+|v|^2\).  Cauchy--Schwarz gives
\(u\cdot v\leq |u|\,|v|\), so the desired squared inequality follows.  Taking
square roots preserves the inequality because both sides are nonnegative.
\end{proof}

\begin{remark*}[Proof structure]
The proof verifies nonnegativity, zero-distance, and symmetry by coordinates.
The triangle inequality is Minkowski's inequality in finite dimensions, proved
from Cauchy--Schwarz.
\end{remark*}

\begin{dependencies}
\begin{itemize}
  \item \hyperref[def:metric-space]{Definition: Metric Space}
  \item \hyperref[def:euclidean-distance-rn]{Definition: Euclidean Distance on \(\mathbb{R}^n\)}
  \item \hyperref[thm:cauchy-schwarz-inequality]{Theorem: Cauchy--Schwarz Inequality}
\end{itemize}
\end{dependencies}

\clearpage
